\makeabstract
    {
    Synthesizing Osteoinductive Graphenic Materials More Efficiently and Cost-Effectively
    }
    {
    Keshav Kumar\textsuperscript{1,2}, Crystal Huynh\textsuperscript{2}, Max (Chenyun) Deng\textsuperscript{2}, Stefanie Sydlik\textsuperscript{2}
    }
    {
    \textsuperscript{1} Department of Chemistry, Amherst College, Amherst, MA\\
\textsuperscript{2} Department of Chemistry, Carnegie Mellon University, Pittsburgh, PA
    }
    {
    Currently, metal implants are frequently used to treat serious bone injuries; however, implants do not match the natural composition or porous structure of healthy bone. Because bone adapts to the loads placed under it, metal implants can absorb the load instead of the bone tissue, leading to bone weakening over time. This research project looks at a novel, more efficient, and cost-effective method for the synthesis of calcium phosphate graphene (CaPG), an osteoinductive graphenic material that can encourage bone regeneration.
    }
    {
    The project examined a novel, aqueous synthetic method for making CaPG intermediates (pE/pK-CDI-GO). Graphene Oxide (GO) was synthesized from graphite using modified Hummer{\textquoteright}s method (using sulfuric acid, sodium nitrate, potassium permanganate, followed by hydrogen peroxide and water), which was then activated with 1,1'-Carbonyldiimidazole (CDI) in water or tetrahydrofuran (THF). Polymer attachment of protected poly-lysine or protected poly-glutamic acid was performed under aqueous conditions. Data on surface functionalization was obtained using X-ray photoelectron spectroscopy (XPS) elemental analysis, with the nitrogen atom percentages measured prior to and following the polymer attachment steps in the aforementioned solvents.
    }
    {
    Nitrogen atom percentages following polymer attachment were lower when compared to the control samples (GO after activation with CDI in the two solvents). For the water control, the nitrogen percentage was {\textasciitilde}2.1\%; following polymer attachment, nitrogen percentage was {\textasciitilde}1.3\% and 0\% for poly-lysine and poly-glutamic acid respectively. For the THF control, the nitrogen percentage was {\textasciitilde}2.0\%; following polymer attachment, nitrogen percentage was {\textasciitilde}0.4\% and {\textasciitilde}1.5\% for poly-lysine and poly-glutamic acid respectively. These findings indicate that the reaction did not proceed as anticipated and significant polymer attachment did not occur.
    }
    {
    While the new synthetic pathway offers theoretical benefits in cost, efficiency, and safety compared to previous methods, the experimental data show that significant polymer attachment was not achieved. Future steps involving deprotected polymers may yield better results, as they have more reactive sites.
    }

    \newpage
    